\section{Schemes, weights, and the limits of the comparison}
\label{sec:framework}

All arithmetic cohomology over the rational integers in this paper lives on
$U_\ell=\operatorname{Spec}\mathbb Z[1/\ell]$. The Tate lattice
$\mathbb Z_\ell(r)$ is the inverse system of finite etale sheaves
$\mu_{\ell^n}^{\otimes r}$, interpreted by continuous cohomology.
The discrete sheaf $\mathbb Q_\ell/\mathbb Z_\ell(-e)$ is its corresponding
divisible torsion coefficient object. These coefficient conventions,
the local conditions, and the compact/discrete duality are part of the
statements; replacing them by rational cohomology loses the finite groups
being bounded.

The weight convention is Deligne's \emph{Weil II}, Definitions 1.2.1--1.2.2,
p.~153, and 1.2.5(iv), p.~154 \cite{DeligneWeilII}: $\mathbb Q_\ell(1)$ has weight $-2$.
Thus the rational Tate object $\mathbb Q_\ell(-e)$ has weight $2e$,
and $\mathbb Q_\ell(e+1)$ has weight $-2(e+1)$. These statements refer
to rational sheaves and geometric Frobenius. They do not assign a
Frobenius purity theorem to a torsion module or prove its vanishing.
The finite $K$-groups at issue have zero rationalization even when their
integral vanishing is unknown. This elementary observation explains why
a weight or sign argument requires the additional integral congruence
and duality statements reconstructed above.

For the archive passage, \emph{SGA 4}, IX.3.2--3.3 supplies the
prime-invertible Kummer sequence and Hilbert 90. The first Chern class
is the Kummer boundary of a line bundle; \emph{SGA 5}, VII,
Proposition 3.4(i), identifies the degree-one term of the total Chern
class with that boundary \cite{SGA5}. The finite root-coordinate gluing proof in
the archive section fixes the precise form needed here. The tame
cover comparison is \emph{SGA 1}, XIII.5.2--5.6, with X.3.6 for the
discrete-valuation Abhyankar lemma. The extension of regular functions
across codimension at least two uses \emph{EGA} $\mathrm{IV}_2$,
5.8.6, 5.9.8--5.9.9 and 5.10.5, as checked below. No general FGA
representability statement supplies the missing arithmetic kernel
or the intrinsic recognition of geometric tangent bases.

The $\lambda$-ring language belongs first to $K_0$: exterior powers
define $\lambda^i$, and the determinant or first gamma-filtration
piece records line bundles. An exact locator is \emph{SGA 6}, V,
Definition 2.1, Example 2.2(b), and Lemma 2.2.1 \cite{SGA6}. Higher $K_{4k}$-groups
require the higher etale comparison used in the next section. Neither
the $K_0$ operations nor a first Chern class is that comparison map.

Crystals and divided powers concern a different coefficient setting,
especially reduction in residue characteristic $\ell$. No new theorem
about a crystal, a divided-power envelope, or a crystalline comparison
is proved in this run. The etale Kummer calculation deliberately removes
$\ell$ from the base; for noninvertible exponents one must use flat
Kummer theory and then state a separate comparison with any crystalline
or prismatic object. The inseparable counterexample in the archive
section proves why the etale assertion itself cannot be transported
unchanged. Assigning one of the present finite Selmer estimates to a
crystalline object without such a map would add an unproved assertion.

Finally, the rational mixed-Tate category over $\operatorname{Spec}
\mathbb Z$ is the category used by Deligne--Goncharov and Brown. Its
base scheme controls ramification; its rational coefficients control
which extensions remain visible. The integral arithmetic lattice of a
pro-$\ell$ image is an additional object. The counterexample at $37$
below shows that these distinctions change the answer to B3.

